\documentclass{beamer}

\usepackage[utf8]{inputenc}
\usepackage[english]{babel}
\usepackage{textpos}
\usepackage{graphicx}
\usepackage{textcomp}
\usepackage{animate}
\usepackage{comment}
%\usepackage{multicol}
%\setlength{\columnsep}{5cm}
\usepackage{amsmath}
\usepackage{amsfonts}
\usepackage{verbatim}
\usepackage{varwidth}
\usepackage{caption}
\usepackage{dcolumn}
\usepackage{amssymb}
\usepackage{algorithm}
\usepackage{algorithmic}
%\usepackage{physics}
\usepackage{hyperref}
\usepackage{scalerel,stackengine}
%\usepackage{fancyvrb}
\usepackage{enumitem}
\setlist[itemize]{label=\textbullet}
\usepackage{tikz}
\usepackage{pgfplots}

%\setbeamertemplate{itemize items}[circle]
%\setbeamertemplate{itemize subitem}[square]

\usefonttheme[onlymath]{serif}

\newcommand{\code}[1]{\texttt{#1}}
\newcommand{\shp}{$>>>$ }
\newcommand{\tab}{\hspace*{22pt}}

\usetheme{Madrid}

\title{Overview of Machine Learning and Predictive Modeling}
%\author{DATA 3550}
\date{}

%\institute{MTSU}
 


%%%%%%%%% BEGIN DOCUMENT %%%%%%%%%
%%%%%%%%% BEGIN DOCUMENT %%%%%%%%%
%%%%%%%%% BEGIN DOCUMENT %%%%%%%%%

\begin{document}

% get title page 
\frame{\titlepage}


\begin{frame}
	\frametitle{Overview of ML}
	
	\begin{alertblock}{What is ML:}
		\footnotesize
		\textbf{Machine learning} is a field of data science/statistics/artificial intelligence which involves creating \textbf{models} to enable computers to ``learn'' patterns in data and make predictions based on these patterns
		\begin{itemize}
			\item A model is a function which can map input data to output data: \\
			$f($input$) = $ output
%			\item Using a dataset of inputs/outputs to learn from, the model can learn the patterns (i.e., find a function $f$) and then, when given new input data, predict its output
%			\item \textbf{Predictive modeling} is a subfield of ML focused on making predictions about the future
		\end{itemize}
		
		\vspace{10pt}
		Example: create a model to predict real estate prices
		\begin{itemize}
			\item Input data: square footage, location, age, number of bedrooms, etc.
			\item Output data: the price (a continuous value)
		\end{itemize}
		
		\vspace{10pt}
		Example: predict whether a patient has cancer
		\begin{itemize}
			\item Input data: age, gender, BMI, family history, environmental factors, etc.
			\item Output data: does/doesn't have cancer (a discrete label)
		\end{itemize}
	\end{alertblock}
\end{frame}

\begin{frame}	
	\frametitle{Overview of ML}
	
	\begin{block}{Some ML terms:}
		\scriptsize
		\begin{itemize}[itemindent=4.0em]
			\item[Dataset:] a collection of inputs/outputs
			\item[Samples:] entries in the dataset
			\item[] (rows in the dataset)
			\item[Features:] the input data (usually $X$, a 2D array $(N_{samp},N_{feat})$)
			\item[] (columns in the dataset)
			\item[] (sometimes called predictors)
			\item[Labels:] the output data (usually $y$, a 1D array $(N_{samp})$) \\
			\item[] (sometimes called ``target'')
			\item[] (a special column in the dataset)
			\item[Model:] a function $f$ which predicts labels based on features
			\item[Training:] the process by which the model ``learns'' the patterns in a dataset
			\item[Testing:] the process by which a model's performance/accuracy is evaluated
			\item[Loss function:] a function which evaluates the amount of error in the model
			\item[] (sometimes called error, must be minimized)
			\item[] (training is minimizing the loss function)
		\end{itemize}
	\end{block}
	
\end{frame}

\begin{frame}	
	\frametitle{Overview of ML}
	
	\begin{block}{A Taxonomy of ML (pt 1):}
		\footnotesize
		One way to classify the different types of ML is based on the nature of the task (i.e., the features and labels):
		\begin{itemize}%[itemindent=1em, align=left]
			\item Supervised learning: a type of learning where the dataset has features AND labels and the goal is to predict the labels based on the features 
			\begin{itemize}
				\item Regression (label is continuous)
				\item Classification (label is categorical/discrete)
			\end{itemize}
			\item Unsupervised learning: a type of learning where the dataset ONLY has features and the goal is to assign a set of reasonable labels
			\begin{itemize}
				\item Clustering
			\end{itemize}
			\item Generative learning:a type of learning that models the distribution of the data so that similar data can be generated (terms features/labels aren't used as much)
			\begin{itemize}
				\item AI chatbots
				\item AI image generators
			\end{itemize}
		\end{itemize}
	\end{block}
	
\end{frame}

\begin{frame}	
	\frametitle{Overview of ML}
	
	\begin{block}{A Taxonomy of ML (pt 2):}
		\footnotesize
		Another way to classify ML methods is based on the nature of the algorithms used:
		\begin{itemize}%[itemindent=1em, align=left]
			\item Linear models: model the relationship between features/labels using linear equations (e.g., linear regression and logistic regression)
			\item Instance-based learning: makes predictions based on the similarity of new instances to existing data points (e.g., $k$-nearest neighbors )
			\item Tree-based methods: use decision trees as the fundamental building block (e.g., decision trees, random forests, and gradient boosting)
			\item Ensemble methods: combine multiple models to improve overall performance (e.g., bagging and boosting)
			\item Deep learning: uses neural networks with multiple layers (e.g., dense NN, convolutional NN, recurrent NN, LSTM's, transformers)
		\end{itemize}
	\end{block}
\end{frame}


\begin{frame}	
	\frametitle{Overview of ML}
	
	\begin{block}{Python libraries for data processing:}
		\begin{itemize}%[itemindent=1em, align=left]
			\item NumPy, SciPy
			\item Pandas
		\end{itemize}
	\end{block}
	
	\begin{block}{Python libraries for visualization:}
		\begin{itemize}%[itemindent=1em, align=left]
			\item Matplotlib, Seaborn
			\item Plotly, Bokeh
		\end{itemize}
	\end{block}
	
	\begin{block}{Python libraries for ML:}
		\begin{itemize}%[itemindent=1em, align=left]
			\item Scikit-learn
			\item Statsmodels
			\item Tensorflow/Keras, PyTorch (neural networks)
			\item XGBoost, LightGBM (gradient boosting)
		\end{itemize}
	\end{block}
	
\end{frame}

\begin{frame}	
	\frametitle{Overview of ML}
	
	\begin{block}{Garbage in, garbage out:}
		No matter how clever and advanced a model is, if it is given bad data...
		\begin{itemize}
			\item useless features that don't realte to the labels
			\item not enough features
			\item missing values, errors in data
			\item insufficient number of samples
			\item too little variety (not all possible cases are covered)
			\item trying to predict beyond the scope of the data (extrapolating)
		\end{itemize}
		then it will make poor predictions
		
%		\vspace{10pt}
%		Thus, \textbf{data preprocessing} is a critical step
	\end{block}
	
\end{frame}

\begin{frame}
	\frametitle{CRISP-DM}
	
	\begin{block}{How to ML:}
		So, what is the correct way to perform ML and (more broadly) data science in general and avoid these and other pitfalls?
		
		\vspace{10pt}
		CRISP-DM: CRoss-Industry Standard Process for Data Mining
		\begin{itemize}
			\item a widely-used methodology for data mining and analytics projects
			\item provides a structured approach to guide the entire DS/ML process
			\item an iterative 6-step process, allowing for revisiting and refining at each step
		\end{itemize}
	\end{block}
	
\end{frame}


\end{document}






